% ----- CHAUPAI 38 -----

\newpage

\subsection*{\deva{अष्टत्रिंशत्तम चौपाई} (Thirty-Eighth Chaupai)}
\addcontentsline{toc}{subsection}{\deva{अष्टत्रिंशत्तम चौपाई} (Thirty-Eighth Chaupai) :  \deva{जो सत...}}

\begin{minipage}[t]{0.55\textwidth}
\vspace{0pt}

{\large \linespread{1.5}\selectfont
\deva{जो सत बार पाठ कर कोई।}\\
\deva{छूटहि बंदि महा सुख होई॥}
\par}

\vspace{0.5em}

{\large \linespread{1.5}\selectfont
\telu{జో సత బార పాఠ కర కోఈ।}\\
\telu{ఛూటహి బంది మహా సుఖ హోఈ॥}
\par}

\vspace{0.5em}

{\large \linespread{1.3}\selectfont
\textit{jo sata bāra pāṭha kara koī |}\\
\textit{chūṭahi bandi mahā sukha hoī ||}
\par}
\end{minipage}%
\hfill
\begin{minipage}[t]{0.42\textwidth}
\vspace{0pt}
\begin{tcolorbox}[anchorbox]
\textbf{The Power of Persistence:} \textit{Sata Baar} (100 times) represents unwavering consistency. In the *Valmiki Ramayana* (*Sundara Kanda, 12:10*), when Hanuman initially fails to find Sita and feels discouraged, he recovers his resolve by declaring: *Anirvedah shriyo mulam* ("Persistence/cheerfulness is the root of all success and prosperity"). He renews his search, proving that mastery and liberation from failure are only built through relentless, repetitive effort.
\vspace{0.5em}\newline Hanuman's ultimate success was built on his absolute refusal to give up, even when his initial efforts yielded no results.\newline\null\hfill\href{https://www.valmiki.iitk.ac.in/sloka?field_kanda_tid=5&language=dv&field_sarga_value=12}{\textit{Sundara Kanda, 12:10}}
\end{tcolorbox}
\end{minipage}

\vspace{0.5em}
\subsubsection*{\textbf{Visual Rhythm Map:}}
\begin{table}[h!]
\centering
\renewcommand{\arraystretch}{1.2}
\setlength{\tabcolsep}{0pt}
\begin{tabularx}{\textwidth}{|*{16}{>{\centering\arraybackslash\hsize=1.0\hsize}X|}}
\hline
\multicolumn{2}{|>{\centering\arraybackslash\hsize=2.0\hsize}X|}{\textbf{\guru{जो}} \par (2)} &
\multicolumn{2}{>{\centering\arraybackslash\hsize=2.0\hsize}X|}{\textbf{\laghu{स}\laghu{त}} \par (2)} &
\multicolumn{3}{>{\centering\arraybackslash\hsize=3.0\hsize}X|}{\textbf{\guru{बा}\laghu{र}} \par (3)} &
\multicolumn{3}{>{\centering\arraybackslash\hsize=3.0\hsize}X|}{\textbf{\guru{पा}\laghu{ठ}} \par (3)} &
\multicolumn{2}{>{\centering\arraybackslash\hsize=2.0\hsize}X|}{\textbf{\laghu{क}\laghu{र}} \par (2)} &
\multicolumn{4}{>{\centering\arraybackslash\hsize=4.0\hsize}X|}{\textbf{\guru{को}\guru{ई}} \par (4)} \\
\hline
\end{tabularx}
\vspace{-1.5em}
\end{table}

\begin{table}[h!]
\centering
\renewcommand{\arraystretch}{1.2}
\setlength{\tabcolsep}{0pt}
\begin{tabularx}{\textwidth}{|*{16}{>{\centering\arraybackslash\hsize=1.0\hsize}X|}}
\hline
\multicolumn{4}{|>{\centering\arraybackslash\hsize=4.0\hsize}X|}{\textbf{\guru{छू}\laghu{ट}\laghu{हि}} \par (4)} &
\multicolumn{3}{>{\centering\arraybackslash\hsize=3.0\hsize}X|}{\textbf{\guru{बं}\laghu{दि}} \par (3)} &
\multicolumn{3}{>{\centering\arraybackslash\hsize=3.0\hsize}X|}{\textbf{\laghu{म}\guru{हा}} \par (3)} &
\multicolumn{2}{>{\centering\arraybackslash\hsize=2.0\hsize}X|}{\textbf{\laghu{सु}\laghu{ख}} \par (2)} &
\multicolumn{4}{>{\centering\arraybackslash\hsize=4.0\hsize}X|}{\textbf{\guru{हो}\guru{ई}} \par (4)} \\
\hline
\end{tabularx}
\vspace{-1.5em}
\end{table}

\vspace{1em}
\textbf{ \deva{सरल-संस्कृतम्}:}\\
 \deva{यः कश्चित् शत-वारम् अस्य पाठं करोति}\\
 \deva{सः बन्धनात् मुच्यते, तस्य महान् सुखं च भवति॥}\\


Whoever recites this one hundred times is freed from all bondage and experiences immense happiness.


\vspace{1em}
\newpage
\textbf{\deva{प्रतिपदार्थः} (Meaning of each word):}
\begin{table}[h!]
\centering
\renewcommand{\arraystretch}{1.2}
\begin{tabularx}{\textwidth}{@{} l l X X @{}}
\toprule
\textbf{Awadhi} & \textbf{Sanskrit} & \textbf{English} & \textbf{Grammar \& Notes} \\
\midrule
\deva{जो कोई} / \telu{జో కోఈ} / \textit{jo koī}  &  \deva{यः कश्चित्}  &  Whoever  &   \\
\deva{सत बार} \gotchaicon / \telu{సత బార} / \textit{sata bāra}  &  \deva{शत-वारम्}  &  100 Times  &  \\ 
\deva{पाठ कर} / \telu{పాఠ కర} / \textit{pāṭha kara}  &  \deva{पाठं करोति}  &  Recites  &   \\
\deva{छूटहि} / \telu{ఛూటహి} / \textit{chūṭahi}  &  \deva{मुच्यते}  &  Is freed / breaks  & Awadhi passive `-hi` \\
\deva{बंदि} / \telu{బంది} / \textit{bandi}  &  \deva{बन्धनात्}  &  From bondage (prison)  &   \\
\deva{महा सुख} / \telu{మహా సుఖ} / \textit{mahā sukha}  &  \deva{महान् सुखम्}  &  Great happiness  &   \\
\deva{होई} / \telu{హోఈ} / \textit{hoī}  &  \deva{भवति}  &  Happens  &   \\
\bottomrule
\end{tabularx}
\end{table}

\begin{tcolorbox}[gotchabox]
\begin{itemize}
    \item \textbf{The \deva{श/ष} $\rightarrow$ \deva{स} Shift:} Yet another perfect example of palatal softening! The strict Sanskrit \deva{शत} (Shata/100) becomes the pure, flowing Awadhi \deva{सत} (Sata).
\end{itemize}
\end{tcolorbox}

\newpage
